{{دیوان
| عنوان = تعلیم اور اس کے نتائج
| ماخذ = ویکی ماخذ
| ماخذ_لنک = https://ur.wikisource.org/wiki/%D8%AA%D8%B9%D9%84%DB%8C%D9%85_%D8%A7%D9%88%D8%B1_%D8%A7%D8%B3_%DA%A9%DB%92_%D9%86%D8%AA%D8%A7%D8%A6%D8%AC
| مقابلہ = اقبال سائبر لائبریری (اقبال اکادمی پاکستان)
| مقابلہ_لنک = https://www.iqbalcyberlibrary.net
}}
=== (تضمین بر شعرِ مُلّا عرشیؔ) ===

<poem>
خوش تو ہیں ہم بھی جوانوں کی ترقّی سے مگر
لبِ خنداں سے نکل جاتی ہے فریاد بھی ساتھ

ہم سمجھتے تھے کہ لائے گی فراغت تعلیم
کیا خبر تھی کہ چلا آئے گا الحاد بھی ساتھ

گھر میں پرویز کے شیریں تو ہوئی جلوہ نُما
لے کے آئی ہے مگر تیشۂ فرہاد بھی ساتھ

“تخمِ دیگر بکف آریم و بکاریم ز نو
کانچہ کشتیم ز خجلت نتواں کرد درو”
</poem>
